\documentclass[11pt]{article}
% Shared preamble for all daily exercises
% Update this file to change settings (padding, parindent, etc.) for all exercises

\usepackage[margin=0.75in]{geometry}
\usepackage{amsmath}
\usepackage{amssymb}

% Custom settings
\setlength{\parindent}{0pt}
\setlength{\parskip}{0.2cm}

\title{Daily Exercise}
\author{Dhruman Gupta}
\date{13th April}

\begin{document}

% Common header for daily exercises
% This file can be included in other LaTeX documents

\maketitle

\noindent
\textbf{Course:} CS-3330, Theory of Computation

\vspace{0.2cm}

\noindent
\textbf{Question:}\noindent 
Show that PATH is logspace-reducible to co-PATH (the problem of deciding whether there is no path from $s$ to $t$).

\textbf{Answer:}

We will show that $PATH \leq_L co\text{-}PATH$.

Let an instance of $PATH$ be a directed graph $G = (V,E)$ with two distinguished vertices $s,t$. We want a logspace computable function $f$ such that
\[
(G,s,t) \in PATH \iff f(G,s,t) \in co\text{-}PATH.
\]

Now, this is not true for an arbitrary directed graph. For example, it is possible that there is a path from $s$ to $t$ and also a path from $t$ to $s$.

So, use the fact that $PATH$ remains complete even when restricted to directed acyclic graphs. Hence, it is enough to give a logspace reduction for DAGs.

Define
\[
f(G,s,t) = (G,t,s).
\]
That is, we keep the graph the same and just swap $s$ and $t$.

Since $G$ is a DAG, it cannot have both a path from $s$ to $t$ and a path from $t$ to $s$, because that would create a directed cycle. Hence, in a DAG,
\[
s \leadsto t \iff t \not\leadsto s.
\]

Therefore,
\[
(G,s,t) \in PATH \iff (G,t,s) \in co\text{-}PATH.
\]

It remains to show that $f$ is computable in logspace. But this is immediate: we output the same encoding of $G$ and swap the encodings of $s$ and $t$, which uses only logarithmic space.

Hence $f$ is a logspace reduction from $PATH$ to $co\text{-}PATH$. Therefore,
\[
PATH \leq_L co\text{-}PATH.
\]

\end{document}
